% !TeX root = ../偏微分方程笔记.tex
% 由主文档通过 \input 载入；请编译主文档。
\setcounter{chapter}{1} % 保留原书章号。
\chapter{积分、收敛与常用不等式}
\label{app:analysis}

\section{勒贝格积分与函数空间}
“几乎处处”表示允许除去一个勒贝格测度为零的集合。$L^p(\Omega)$ 的元素是几乎处处相等的可测函数的等价类，范数为
\[
 \|f\|_p=\left(\int_\Omega|f|^p\,dx\right)^{1/p}\quad(1\le p<\infty),\qquad
 \|f\|_\infty=\operatorname*{ess\,sup}_{x\in\Omega}|f(x)|.
\]
本性上确界是忽略零测集后成立的最小本质上界。局部可积 $L^1_{\mathrm{loc}}$ 要求在每个紧子集上可积，不要求全区域积分有限。

若 $|\Omega|<\infty$ 且 $q\ge p$，赫尔德不等式给出 $\|f\|_p\le|\Omega|^{1/p-1/q}\|f\|_q$。在无限测度区域上不能任意沿用这个包含关系。范数收敛也不自动给出原序列逐点收敛，因此初值与级数解必须标明收敛空间。

\section{交换运算与收敛定理}
\begin{description}
 \item[单调收敛定理] 若可测函数 $0\le f_j\uparrow f$，则 $\int f_j\uparrow\int f$，允许积分为无穷。
 \item[控制收敛定理] 若 $f_j\to f$ 几乎处处，且存在同一个 $g\in L^1$ 使 $|f_j|\le g$，则 $\int f_j\to\int f$。控制函数不能随 $j$ 无限制变化。
 \item[法图引理] 若 $f_j\ge0$，则 $\int\liminf_j f_j\le\liminf_j\int f_j$。这类下半连续关系常用于极小化序列。
 \item[托内利与富比尼定理] 在本书使用的勒贝格乘积空间上，非负可测函数可由托内利定理交换积分次序；绝对可积函数可由富比尼定理交换积分次序。非绝对可积的带符号积分需要另作判断。
\end{description}

积分号下求导的一组充分条件是：参数附近导数存在，导数被一个与参数无关的可积函数控制，且原函数在一个参数值处可积。此时可把差商写成导数积分并应用控制收敛。级数逐项求导的一组常用充分条件是导数级数一致收敛且原级数在一点收敛；若还要更高阶导数，就必须对相应阶数重复核查条件。

\section{估计与格朗沃尔不等式}
若 $1<p,q<\infty$、$1/p+1/q=1$，\emph{赫尔德不等式（H\"older inequality）}为 $\int|fg|\le\|f\|_p\|g\|_q$；$p=q=2$ 即\emph{柯西--施瓦茨不等式（Cauchy--Schwarz inequality）}。端点 $p=1,q=\infty$ 也成立。

能量估计常用参数化\emph{杨氏不等式（Young's inequality）}
\[
 ab\le\frac\varepsilon2 a^2+\frac1{2\varepsilon}b^2,\qquad a,b\ge0,\quad\varepsilon>0,
\]
它由 $(\sqrt\varepsilon a-b/\sqrt\varepsilon)^2\ge0$ 得到。选择 $\varepsilon$ 的目的是让第一项被方程左侧的正能量吸收，代价是另一常数增大。

\begin{lemma}[格朗沃尔不等式（Gronwall's inequality）]
设连续函数 $y\ge0$、常数 $a\ge0$ 和非负可积函数 $b$ 满足
\[
 y(t)\le a+\int_0^t b(s)y(s)ds,
\]
则 $y(t)\le a\exp(\int_0^tb(s)ds)$。
\end{lemma}
\begin{proof}
先设 $a>0$，令 $Y(t)=a+\int_0^tby$。则 $y\le Y$ 且 $Y'\le bY$ 几乎处处成立，乘以 $\exp(-\int_0^tb)$ 后导数非正，便得结论。$a=0$ 时先把右端加上任意 $\varepsilon>0$，再令 $\varepsilon\downarrow0$。
\end{proof}
相应地，若绝对连续的 $y$ 满足 $y'+\lambda y\le r(t)$，则乘以积分因子得 $y(t)\le e^{-\lambda t}y(0)+\int_0^te^{-\lambda(t-s)}r(s)ds$。每次使用估计时，应说明常数依赖哪些参数，特别要指出是否独立于近似维数或网格步长。

